Read in Gauge Theory Gravity: the metric $g$ below is the derived metric of the gravitational field $(\hb, \Omega)$ on flat spacetime, the manifold is the invertible region, and the tractor connection is the flat-space $\Cl(2,4)$ gauge field of Theorem \claimref{synthesis.gtg-tractor-gauge-field-reproduces-tractor-connection}; the index computation below is its reading in the splitting of $g$ (Fact \claimref{synthesis.imported-gtg-dictionary}).

Substitute $(\sigma, \mu, \rho) = (1, 0, 0)$ into \eqref{eq:tractorconn}: the first slot is $\nabla_a 1 - 0 = 0$; the second is $0 + g_{ab}\cdot 0 - \Sch_{ab}\cdot 1 = -\Sch_{ab}$; the third is $0 + \Sch_a{}^b \cdot 0 = 0$. So $\nabla I_g = 0$ iff $\Sch_{ab} = 0$ iff $\Ric = 0$ by Lemma \claimref{interior.metric-parallelism-trace-schouten-condition}(ii), and Fact \claimref{interior.imported-equivalence-einstein-equations} converts Ricci-flatness into the GTG field equations on the region of invertibility. The norm is $Q(I_g) = g(0,0) + 2 \cdot 1 \cdot 0 = 0$. Covariance is the statement that both sides of the equivalence are conditions on a section of an invariantly defined bundle with its canonical connection (Fact \claimref{interior.imported-tractor-connection}), the section itself determined by the given scale $g$: the splitting used to write the proof is scale-dependent; the bundle, the connection, and the parallelism condition on the given section are not.
